% Part: computability
% Chapter: computability-theory
% Section: complement-ce

\documentclass[../../../include/open-logic-section]{subfiles}

\begin{document}

\olfileid{cmp}{thy}{cmp}
\olsection{المجموعات القابلة للتعداد حسابيًا ليست منغلقة تحت المتممة}

إذا كانت $A$ قابلة للتعداد حسابيًا، أيلزم أن تكون متممة~$A$، أعني
$\Complement{A}=\Nat\setminus A$، قابلة للتعداد حسابيًا؟
ليس ذلك بلازم، كما تبيّن المبرهنة التالية ونتيجتها.

\begin{thm}
\ollabel{thm:ce-comp}
لكل مجموعة أعداد طبيعية $A$، تتكافأ قابلية $A$ للحساب
مع كون $A$ و$\Complement{A}$ جميعًا قابلتين للتعداد
حسابيًا.
\end{thm}

\begin{proof}
أما الوجه الأول فظاهر؛ فإن قابلية $A$ للحساب تستلزم قابلية $\Complement{A}$
له ($\Char{A}=1\tsub\Char{\Complement{A}}$)، ومن ثم
تكونان كلتاهما قابلتين للتعداد حسابيًا.

وأما العكس، فلنفرض أن $A$ و~$\Complement{A}$ قابلتان
للتعداد حسابيًا. نجعل $A$ مجال~$\cfind{d}$،
و$\Complement{A}$ مجال~$\cfind{e}$، ونعرّف $h$ كما يأتي:
\[
h(x) = \umin{s}{(T(d,x,s) \lor T(e,x,s))}.
\]
فالدالة~$h$ تبحث عند الدخل~$x$ عن سجل حساب متوقف لـ
~$\cfind{d}$ أو عن سجل حساب متوقف لـ~$\cfind{e}$. فإذا كان $x\in A$،
% تصحيح صياغة كلاسيكية (C242-01): وُحّد حرف الجر في معمولي البحث المتوازيين.
نجحت في الحالة الأولى، وإذا كان $x\in\Complement{A}$، نجحت في الحالة
الثانية. فلا بد من نجاح البحث؛ وتكون~$h$ كلية قابلة للحساب. ثم إن كل
~$x$ يحقق $x\in A$ إذا وفقط إذا صدقت $T(d,x,h(x))$، أي إذا كانت
$\cfind{d}$ هي المعرفة. وحيث تقبل العلاقة $T(d,x,h(x))$ الحساب،
تقبل~$A$ الحساب أيضًا.
\end{proof}

\begin{explain}
ويقرب ذلك بوصف إجرائي: لقرار $A$ نبحث عند الدخل
$x$ عن سجل توقف لـ$\cfind{d}$ أو $\cfind{e}$. ولا يفوتنا
أحدهما؛ فإن توقف $\cfind{d}$، كان $x$ من~$A$، وإلا كان $x$ من
~$\Complement{A}$.
\end{explain}

\begin{cor}
\ollabel{cor:comp-k}
$\Complement{K_0}$ غير قابلة للتعداد حسابيًا.
\end{cor}

\begin{proof}
قد علمنا قابلية $K_0$ للتعداد حسابيًا مع عدم قابليتها للحساب. ولو كانت
$\Complement{K_0}$ قابلة للتعداد حسابيًا، لكانت $K_0$ قابلة للحساب
بحسب \olref{thm:ce-comp}، خلافًا لـ\olref[ncp]{thm:K-0}.
\end{proof}

\end{document}
